\chapter{Native 细节全拆：炼星、网格、环缓冲、两套判决}
\label{ch:native-details}

\section{炼星 Embed：四步药剂}

\begin{metaphor}[炼金只看眼前四味药]
位置 $p$ 的星号只由 \texttt{seed} 与 $b_{p-3..p}$ 决定。\\
\textbf{不}把从文件头累加的药水端进来 —— 否则挪 1 字节，整条河的星号全漂。
\end{metaphor}

\begin{figure}[htbp]
\centering
\begin{tikzpicture}[font=\small]
  \foreach \i/\lab in {0/p-3,1/p-2,2/p-1,3/p} {
    \node[bytecell, fill=Gen51Color!25] (b\i) at (\i*1.0,0) {\lab};
  }
  \node[softbox=Gen51Color, below=0.8cm of b1] (y) {$Y$};
  \draw[-{Stealth}, thick] (b0)--(y);
  \draw[-{Stealth}, thick] (b1)--(y);
  \draw[-{Stealth}, thick] (b2)--(y);
  \draw[-{Stealth}, thick] (b3)--(y);
\end{tikzpicture}
\caption{闭合多项式与四步 LCG 迭代必须同一杯药（bit-identical）。}
\label{fig:four-herbs}
\end{figure}

真名：\texttt{LandmarkEmbedY4}；$m=1664525$，$c=1013904223$。

\section{网格落点：二的幂掩码}

\texttt{event\_stride} 先收成最近 $2^\ell$（平局偏小）。然后
\begin{equation}
Y \mathbin{\&}(S-1)=0 \quad\Leftrightarrow\quad Y\equiv 0\pmod S.
\end{equation}

\begin{toybox}[平局偏小]
想要 100：落在 64 与 128 正中？规则选 64。\\
避免 ceil 一把拉到 128 导致突然``天空太稀、块巨大''。
\end{toybox}

\section{时间轴 $T$：本块相对尺}

$T$ 用块内相对偏移，不是全文件绝对尺。切后星图清空，下一块从 0 再数。\\
坐标再对 \texttt{q\_mod}（默认 1024）取模，避免数轴无限涨。

\section{环缓冲：最旧在 \texttt{ring\_start}}

\begin{figure}[htbp]
\centering
\begin{tikzpicture}[font=\small]
  \foreach \i/\ang in {0/90,1/18,2/-54,3/-126,4/162} {
    \fill[Gen51Color!40] ({1.6*cos(\ang)},{1.6*sin(\ang)}) circle (0.2);
    \node[font=\tiny] at ({1.6*cos(\ang)},{1.6*sin(\ang)}) {\i};
  }
  \draw[-{Stealth}, thick, CutRed] (0,0)--({1.0*cos(90)},{1.0*sin(90)});
  \node[font=\footnotesize, CutRed] at (0,-2.3) {ring\_start 指向最旧；满则覆盖并前移};
\end{tikzpicture}
\caption{满窗时逻辑顺序 $=$ 从 ring\_start 顺时针数 $k$ 颗。}
\label{fig:ring}
\end{figure}

\section{Tri 判决}

$|\mathrm{Flip}_t-\mathrm{Flip}_{t-1}|\ge\tau$，且距上次切 $\ge\min$（\texttt{MinGapOk}）。\\
$\Delta\mathrm{Flip}$ 常是 0/1，所以\textbf{真正拧松紧的是 stride 网格}，不是把 $\tau$ 拧到 100。

\section{PH 判决：翻面 $\land$ 连通（$\land$ 寿命）}

\begin{enumerate}[nosep]
  \item 钉星前拍一张三档 $\beta0$；
  \item 钉星后再拍；
  \item 任一档变 $\Rightarrow$ ph0OK；
  \item 同时要 flipOK；
  \item 若开 persist：$\max\beta0-\min\beta0\ge\delta$。
\end{enumerate}

\begin{figure}[htbp]
\centering
\begin{tikzpicture}[font=\small]
  \node[softbox=Gen51Color] (a) {flipOK};
  \node[softbox=Gen51Color, right=0.5cm of a] (b) {ph0OK};
  \node[softbox=Gen51Color, right=0.5cm of b] (c) {persist?};
  \node[softbox=CutRed, right=0.5cm of c] (d) {切};
  \draw[-{Stealth}, thick] (a)--(b)--(c)--(d);
\end{tikzpicture}
\caption{与 Tri 必须刀口不同（门禁 Tri$\neq$PH）。}
\label{fig:phn-and}
\end{figure}

嵌套 $\{1,4,16\}$：边按距离排序一次，从小半径扫到大 —— 像一次过安检三道门。

\section{切后拍掉星空}

\begin{pitfall}[不拍掉]
下一块继承``半截天空''，1 字节移位后无法重新锁边，$\mathrm{reuse}$ 崩盘。\\
真名：\texttt{ResetLandmarkWindow}。
\end{pitfall}

\section{定理剪枝（人话）}

\begin{center}
\begin{tabular}{@{}lp{6.5cm}@{}}
\toprule
留下 & 扔掉 \\
\midrule
小窗 $H_0$、定点 Flip、寿命跨度 & 全球 PH、在线 bottleneck、$H_1+$、神经定理 \\
reuse 评测可讲稳定性故事 & 热路径不算 $W_\infty$ \\
\bottomrule
\end{tabular}
\end{center}
